\subsection{M-平坦模}

定义 1. 设 A 是环，E 是右 A-模，M 是左 A-模。称 E 对 M 是平坦的（或 M-平坦的），如果对每个左 A-模 M' 和每个单同态 \(v: M' \to M\) ，同态 \(1_{E} \otimes v\) ： E \(\otimes_{A} M' \to E \otimes_{A} M\) 是单射。

对任意右 A-模 N，N-平坦左模的概念可类似地定义。说一个右 A-模 E 对一个左 A-模 M 是平坦的，等价于说 E 视为左 \(A^{0}\) -模（回顾 \(A^{0}\) 表示 A 的反环）时对右 \(A^{0}\) -模 M 是平坦的。

引理 3. 为使一个右 A-模 E 是 M-平坦的，必要且充分的是：对 M 的每个有限生成子模 M'，典范同态

\[
\mathbf {1} _ {\mathbf {E}} \otimes j: \mathbf {E} \otimes_ {\mathbf {A}} \mathbf {M} ^ {\prime} \rightarrow \mathbf {E} \otimes_ {\mathbf {A}} \mathbf {M}
\]

（j 为典范单射 \(\mathbf{M}'\to \mathbf{M}\) ）是单射的。

假设这一条件成立，并设 N 是 M 的任意子模。设在 \(E \otimes_{A} M\) 中，一个元素

\[
z = \sum_ {i} x _ {i} \otimes y _ {i} \in \mathbf {E} \otimes_ {\mathbf {A}} \mathbf {N} (x _ {i} \in \mathbf {E}, y _ {i} \in \mathbf {N})
\]

的典范像为零，并设 \(M'\) 是由诸 \(y_{i}\) 生成的 N 的有限生成子模；由于由假设复合映射 \(E \otimes_{A} M' \to E \otimes_{A} N \to E \otimes_{A} M\) 是单射的，把和 \(z' = \sum_{i} x_{i} \otimes y_{i}\) 视为 \(E \otimes_{A} M'\) 的元素时为零。由于 z 是 \(z'\) 的像，我们也有 z = 0，引理得证。

引理 4. 设 E 是右 A-模，M 是左 A-模，且 E 是 M-平坦的。若 N 是 M 的子模或商模，则 E 是 N-平坦的。

N 是子模的情形是容易的，因为若 \(N'\) 是 N 的子模，则复合同态

\[
\mathrm{E} \otimes_ {\mathrm{A}} \mathrm{N} ^ {\prime} \rightarrow \mathrm{E} \otimes_ {\mathrm{A}} \mathrm{N} \rightarrow \mathrm{E} \otimes_ {\mathrm{A}} \mathrm{M}
\]

是单射的，因此 \(\mathbf{E} \otimes_{\mathbf{A}} \mathbf{N}' \to \mathbf{E} \otimes_{\mathbf{A}} \mathbf{N}\) 也是单射的。于是设 \(\mathbf{N}\) 是 \(\mathbf{M}\) 的一个商模，即存在正合列 \(0 \to \mathbf{R} \xrightarrow{i} \mathbf{M} \xrightarrow{p} \mathbf{N} \to 0\) 。设 \(\mathbf{N}'\) 是 \(\mathbf{N}\) 的一个子模，并令 \(\mathbf{M}' = \bar{p}^{-1}(\mathbf{N}')\) 。用 \(i'\) 表示 \(\mathbf{R}\) 到 \(\mathbf{M}'\) 的、与 \(i, p'\) 中 i 有相同图的映射，以及 \(\mathbf{M}' \to \mathbf{N}'\) 的、与 \(p\) 在 \(\mathbf{M}'\) 上的限制有相同图的满射，用 \(r\) 表示 \(\mathbf{R}\) 到 \(\mathbf{R}\) 的恒等映射， \(m\) 表示典范单射 \(\mathbf{M}' \to \mathbf{M}\) ， \(n\) 表示典范单射 \(\mathbf{N}' \to \mathbf{N}\) 。图

\[
\begin{array}{c} 0 \longrightarrow \mathrm{R} \xrightarrow {i ^ {\prime}} \mathrm{M} ^ {\prime} \xrightarrow {p ^ {\prime}} \mathrm{N} ^ {\prime} \longrightarrow 0 \\ r \Biggl \downarrow \qquad m \Biggl \downarrow \qquad n \Biggl \downarrow \\ 0 \longrightarrow \mathrm{R} \xrightarrow [ i ]{} \mathrm{M} \xrightarrow [ p ]{} \mathrm{N} \longrightarrow 0 \end{array}
\]

是交换的且它的行是正合的。

为书写简单，我们令 \(\mathrm{T}(\mathbf{Q}) = \mathrm{E} \otimes_{\mathbf{A}} \mathbf{Q}\) （对每个左 A-模 \(\mathbf{Q}\) ），令 \(\mathrm{T}(v) = 1_{\mathbf{E}} \otimes v\) （对左 A-模的每个同态 \(v\) ）。图

\[
\begin{array}{l} \mathrm{T} (\mathrm{R}) \xrightarrow {\mathrm{T} (i ^ {\prime})} \mathrm{T} (\mathrm{M} ^ {\prime}) \xrightarrow {\mathrm{T} (p ^ {\prime})} \mathrm{T} (\mathrm{N} ^ {\prime}) \longrightarrow 0 \\ \mathrm{T} (r) \Bigg \downarrow \qquad \qquad \qquad \qquad \qquad \qquad \qquad \qquad \qquad \qquad \qquad \qquad \qquad \qquad \qquad \qquad \qquad \qquad \qquad \qquad \qquad \qquad \qquad \qquad \qquad \qquad \qquad \qquad \qquad \qquad \qquad \qquad \qquad \qquad \\ \mathrm{T} (\mathrm{R}) \xrightarrow [ \mathrm{T} (i) ]{} \mathrm{T} (\mathrm{M}) \xrightarrow [ \mathrm{T} (p) ]{} \mathrm{T} (\mathrm{N}) \longrightarrow 0 \end{array}
\]

是交换的，且由第 1 目引理 1，它的行是正合的。此外，由于 E 是 M-平坦的，同态 T(m) 是单射。由于 T(r) 和 T(p') 是满射，由 §1 第 4 目命题 2 的推论 2 得 T(n) 是单射，这就证明了引理。

引理 5. 设 \((\mathbf{M}_{\iota})_{\iota \in \mathbf{I}}\) 是一族左 A-模， \(\mathbf{M} = \bigoplus_{\iota \in \mathbf{I}} \mathbf{M}_{\iota}\) 是它们的直和，E 是一个右 A-模。若对所有 \(\iota \in \mathbf{I}\) ， E 对 \(\mathbf{M}_{\iota}\) 都是平坦的，则 E 对 M 是平坦的。

\begin{enumerate}
\def\labelenumi{(\alph{enumi})}
\tightlist
\item
  先设 \(I = \{1, 2\}\) ，并设 \(M'\) 是 \(M = M_{1} \oplus M_{2}\) 的子模，其中 \(M_{1}\) 和 \(M_{2}\) 典范地等同于 M 的子模。用 \(M_{1}'\) 表示交 \(M' \cap M_{1}\) ，用 \(M_{2}'\) 表示 \(M'\) 在 \(M_{2}\) 中、经 M 到 \(M_{2}\) 上的典范投影 p 之下的像。我们有一个图
\end{enumerate}

\[
\begin{array}{c} 0 \longrightarrow \mathrm{M} _ {1} ^ {\prime} \xrightarrow {i ^ {\prime}} \mathrm{M} ^ {\prime} \xrightarrow {p ^ {\prime}} \mathrm{M} _ {2} ^ {\prime} \longrightarrow 0 \\ v _ {1} \Biggl \downarrow \qquad \qquad \qquad \qquad \qquad \qquad \qquad \qquad \qquad \qquad \qquad \qquad \qquad v _ {2} \Biggl \downarrow \\ 0 \longrightarrow \mathrm{M} _ {1} \xrightarrow [ i ]{} \mathrm{M} \xrightarrow [ p ]{} \mathrm{M} _ {2} \longrightarrow 0 \end{array}
\]

其中 \(v_{1}, v, v_{2}, i, i'\) 是典范单射，而 \(p'\) 是与 p 在 M' 上的限制有相同图的映射，它是满射。立即可验证这个图是交换的且它的行是正合的。以证明引理 4 时相同的含义使用 T(Q) 和 T(v) ，我们有一个交换图

\[
\begin{array}{c} \mathrm{T} (\mathrm{M} _ {1} ^ {\prime}) \xrightarrow {\mathrm{T} (i ^ {\prime})} \mathrm{T} (\mathrm{M} ^ {\prime}) \xrightarrow {\mathrm{T} (p ^ {\prime})} \mathrm{T} (\mathrm{M} _ {2} ^ {\prime}) \\ \mathrm{T} (v _ {1}) \Biggl \downarrow \qquad \qquad \qquad \qquad \qquad \qquad \qquad \qquad \qquad \qquad \qquad \qquad \qquad \qquad \qquad \qquad \qquad \qquad \qquad \qquad \qquad \qquad \qquad \qquad \qquad \qquad \qquad \qquad \qquad \qquad \qquad \qquad \qquad \qquad \\ \mathrm{T} (\mathrm{M} _ {1}) \xrightarrow [ \mathrm{T} (i) ]{} \mathrm{T} (\mathrm{M}) \xrightarrow [ \mathrm{T} (p) ]{} \mathrm{T} (\mathrm{M} _ {2}) \end{array}
\]

由第 1 目的引理 1，这个图的两行都是正合的；由于 E 对 \(M_{1}\) 和 \(M_{2}\) 是平坦的， \(T(v_{1})\) 和 \(T(v_{2})\) 是单射；此外，由第 1 目的引理 2， \(T(i)\) 是单射。于是 §1 第 4 目命题 2 的推论 2 表明 \(T(v)\) 是单射，因而 E 是 M-平坦的。

\begin{enumerate}
\def\labelenumi{(\alph{enumi})}
\setcounter{enumi}{1}
\item
  若 I 是有 \(n\) 个元素的有限集，则利用 (a) 对 \(n\) 作归纳即得引理。
\item
  在一般情形，设 \(\mathbf{M}'\) 是 \(\mathbf{M}\) 的一个有限生成子模。则存在一个有限子集 \(J\) ，它包含在指标集 \(I\) 中，使得 \(\mathbf{M}'\) 包含于直和 \(\mathbf{M}_J = \bigoplus_{i \in J} \mathbf{M}_i\) 。由 (b) ， \(E\) 对 \(\mathbf{M}_J\) 是平坦的；因此典范同态 \(\mathbf{E} \otimes_{\mathbf{A}} \mathbf{M}' \to \mathbf{E} \otimes_{\mathbf{A}} \mathbf{M}_J\) 是单射的。另一方面，由于 \(\mathbf{M}_J\) 是 \(\mathbf{M}\) 的一个直因子，典范同态 \(\mathbf{E} \otimes_{\mathbf{A}} \mathbf{M}_J \to \mathbf{E} \otimes_{\mathbf{A}} \mathbf{M}\) 是单射的（第 1 目，引理 2）。取其复合，即得 \(\mathbf{E} \otimes_{\mathbf{A}} \mathbf{M}_J \to \mathbf{E} \otimes_{\mathbf{A}} \mathbf{M}\) 是单射的，从而由引理 3 ， \(E\) 对 \(M\) 是平坦的。
\end{enumerate}

